% Options for packages loaded elsewhere
\PassOptionsToPackage{unicode}{hyperref}
\PassOptionsToPackage{hyphens}{url}
%
\documentclass[
  11pt,
]{article}
\usepackage{amsmath,amssymb}
\usepackage{iftex}
\ifPDFTeX
  \usepackage[T1]{fontenc}
  \usepackage[utf8]{inputenc}
  \usepackage{textcomp} % provide euro and other symbols
\else % if luatex or xetex
  \usepackage{unicode-math} % this also loads fontspec
  \defaultfontfeatures{Scale=MatchLowercase}
  \defaultfontfeatures[\rmfamily]{Ligatures=TeX,Scale=1}
\fi
\usepackage{lmodern}
\ifPDFTeX\else
  % xetex/luatex font selection
\fi
% Use upquote if available, for straight quotes in verbatim environments
\IfFileExists{upquote.sty}{\usepackage{upquote}}{}
\IfFileExists{microtype.sty}{% use microtype if available
  \usepackage[]{microtype}
  \UseMicrotypeSet[protrusion]{basicmath} % disable protrusion for tt fonts
}{}
\makeatletter
\@ifundefined{KOMAClassName}{% if non-KOMA class
  \IfFileExists{parskip.sty}{%
    \usepackage{parskip}
  }{% else
    \setlength{\parindent}{0pt}
    \setlength{\parskip}{6pt plus 2pt minus 1pt}}
}{% if KOMA class
  \KOMAoptions{parskip=half}}
\makeatother
\usepackage{xcolor}
\usepackage[margin=2.2cm]{geometry}
\setlength{\emergencystretch}{3em} % prevent overfull lines
\providecommand{\tightlist}{%
  \setlength{\itemsep}{0pt}\setlength{\parskip}{0pt}}
\setcounter{secnumdepth}{-\maxdimen} % remove section numbering
\ifLuaTeX
  \usepackage{selnolig}  % disable illegal ligatures
\fi
\IfFileExists{bookmark.sty}{\usepackage{bookmark}}{\usepackage{hyperref}}
\IfFileExists{xurl.sty}{\usepackage{xurl}}{} % add URL line breaks if available
\urlstyle{same}
\hypersetup{
  pdftitle={Worksheet 1A: Fractions Fundamentals},
  hidelinks,
  pdfcreator={LaTeX via pandoc}}

\title{Worksheet 1A: Fractions Fundamentals}
\author{}
\date{}

\begin{document}
\maketitle

\hypertarget{part-a-numerical-fractions-12-problems}{%
\subsection{Part A: Numerical Fractions (12
problems)}\label{part-a-numerical-fractions-12-problems}}

Simplify or evaluate. Show your work.

\textbf{1.} \(\dfrac{3}{4} + \dfrac{5}{6}\)

\textbf{2.} \(\dfrac{7}{8} - \dfrac{2}{3}\)

\textbf{3.} \(\dfrac{5}{9} \times \dfrac{3}{10}\)

\textbf{4.} \(\dfrac{4}{5} \div \dfrac{2}{15}\)

\textbf{5.} \(\dfrac{2}{3} + \dfrac{1}{4} - \dfrac{1}{6}\)

\textbf{6.} \(\dfrac{7}{12} \times \dfrac{8}{21} \times \dfrac{3}{2}\)

\textbf{7.} \(\dfrac{\frac{3}{4}}{\frac{9}{16}}\)

\textbf{8.} \(\dfrac{\frac{2}{5} + \frac{1}{3}}{\frac{1}{2}}\)

\textbf{9.} \(2\dfrac{1}{3} \times 1\dfrac{1}{2}\)

\textbf{10.} \(4\dfrac{2}{5} \div 1\dfrac{1}{10}\)

\textbf{11.} \(\dfrac{5!}{4!}\)

\textbf{12.} \(\dfrac{7!}{5! \cdot 2!}\)

\begin{center}\rule{0.5\linewidth}{0.5pt}\end{center}

\hypertarget{part-b-pattern-sums-8-problems}{%
\subsection{Part B: Pattern Sums (8
problems)}\label{part-b-pattern-sums-8-problems}}

Evaluate each sum. Look for patterns.

\textbf{13.} \(1 + \dfrac{1}{2} + \dfrac{1}{4}\)

\textbf{14.} \(1 + \dfrac{1}{2} + \dfrac{1}{4} + \dfrac{1}{8}\)

\textbf{15.}
\(1 + \dfrac{1}{2} + \dfrac{1}{4} + \dfrac{1}{8} + \dfrac{1}{16}\)

\textbf{16.} Based on problems 13--15, what value does this sum approach
if we keep adding terms forever?

\textbf{17.}
\(\dfrac{1}{1 \cdot 2} + \dfrac{1}{2 \cdot 3} + \dfrac{1}{3 \cdot 4}\)

\emph{(Hint: Try writing each fraction as a difference of two simpler
fractions)}

\textbf{18.}
\(\dfrac{1}{1 \cdot 2} + \dfrac{1}{2 \cdot 3} + \dfrac{1}{3 \cdot 4} + \dfrac{1}{4 \cdot 5}\)

\textbf{19.} Based on 17--18, if we kept adding terms
\(\dfrac{1}{n(n+1)}\) forever, what would the sum approach?

\textbf{20.}
\(1 + \dfrac{1}{1!} + \dfrac{1}{2!} + \dfrac{1}{3!} + \dfrac{1}{4!}\)

\emph{(This sum, extended infinitely, equals \(e \approx 2.718...\))}

\begin{center}\rule{0.5\linewidth}{0.5pt}\end{center}

\hypertarget{part-c-algebraic-fractions-simplification-10-problems}{%
\subsection{Part C: Algebraic Fractions --- Simplification (10
problems)}\label{part-c-algebraic-fractions-simplification-10-problems}}

Simplify completely. Assume no denominators equal zero.

\textbf{21.} \(\dfrac{6p^2}{9p}\)

\textbf{22.} \(\dfrac{15m^3n^2}{25mn^4}\)

\textbf{23.} \(\dfrac{t^2 - 4}{t + 2}\)

\textbf{24.} \(\dfrac{r^2 - 9}{r^2 + 6r + 9}\)

\textbf{25.} \(\dfrac{2k^2 + 4k}{k^2 + 2k}\)

\textbf{26.} \(\dfrac{h^2 - 5h + 6}{h^2 - 4}\)

\textbf{27.} \(\dfrac{n!}{(n-1)!}\)

\textbf{28.} \(\dfrac{(n+1)!}{n!}\)

\textbf{29.} \(\dfrac{(n+2)!}{n! \cdot 2!}\)

\textbf{30.} \(\dfrac{6s^2 - s - 2}{2s^2 + 3s - 2}\)

\begin{center}\rule{0.5\linewidth}{0.5pt}\end{center}

\hypertarget{part-d-algebraic-fractions-operations-10-problems}{%
\subsection{Part D: Algebraic Fractions --- Operations (10
problems)}\label{part-d-algebraic-fractions-operations-10-problems}}

Perform the indicated operations and simplify.

\textbf{31.} \(\dfrac{2}{p} + \dfrac{3}{q}\)

\textbf{32.} \(\dfrac{5}{t+1} - \dfrac{2}{t-1}\)

\textbf{33.} \(\dfrac{3}{m} + \dfrac{2}{n} - \dfrac{1}{mn}\)

\textbf{34.} \(\dfrac{r}{r+2} + \dfrac{2}{r-2}\)

\textbf{35.} \(\dfrac{1}{k-1} - \dfrac{1}{k+1}\)

\textbf{36.} \(\dfrac{2s}{s+3} \times \dfrac{s^2 - 9}{4s^2}\)

\textbf{37.} \(\dfrac{w^2 - 4}{w + 5} \div \dfrac{w - 2}{w^2 + 5w}\)

\textbf{38.} \(\dfrac{1}{u} + \dfrac{1}{u+1} + \dfrac{1}{u+2}\)

\textbf{39.}
\(\dfrac{p^2 - 1}{p^2 + 2p} \cdot \dfrac{p^2 + 4p + 4}{p^2 + p - 2} \div \dfrac{p + 2}{p}\)

\textbf{40.} If \(\dfrac{1}{m} + \dfrac{1}{n} = \dfrac{1}{k}\), solve
for \(k\) in terms of \(m\) and \(n\).

\begin{center}\rule{0.5\linewidth}{0.5pt}\end{center}

\hypertarget{part-e-complex-fractions-6-problems}{%
\subsection{Part E: Complex Fractions (6
problems)}\label{part-e-complex-fractions-6-problems}}

Simplify each expression.

\textbf{41.}
\(\dfrac{\frac{1}{p} + \frac{1}{q}}{\frac{1}{p} - \frac{1}{q}}\)

\textbf{42.} \(\dfrac{\frac{2}{t+1} - \frac{1}{t}}{\frac{3}{t}}\)

\textbf{43.} \(\dfrac{1 - \frac{1}{n^2}}{1 + \frac{1}{n}}\)

\textbf{44.} \(\dfrac{\frac{m}{n} + \frac{n}{m}}{\frac{m^2 - n^2}{mn}}\)

\textbf{45.} \(\dfrac{1}{1 + \frac{1}{1 + \frac{1}{k}}}\)

\textbf{46.} \(\dfrac{\frac{1}{h} - \frac{1}{h+1}}{h}\)

\begin{center}\rule{0.5\linewidth}{0.5pt}\end{center}

\hypertarget{part-f-repeated-multiplication-shrinking-vs-growing-6-problems}{%
\subsection{Part F: Repeated Multiplication --- Shrinking vs Growing (6
problems)}\label{part-f-repeated-multiplication-shrinking-vs-growing-6-problems}}

\textbf{47.} Calculate each: - (a) \(\dfrac{3}{4}\) - (b)
\(\left(\dfrac{3}{4}\right)^2\) - (c) \(\left(\dfrac{3}{4}\right)^3\) -
(d) \(\left(\dfrac{3}{4}\right)^4\) - (e) What happens to
\(\left(\dfrac{3}{4}\right)^n\) as \(n\) gets very large?

\textbf{48.} Calculate each: - (a) \(\dfrac{5}{4}\) - (b)
\(\left(\dfrac{5}{4}\right)^2\) - (c) \(\left(\dfrac{5}{4}\right)^3\) -
(d) \(\left(\dfrac{5}{4}\right)^4\) - (e) What happens to
\(\left(\dfrac{5}{4}\right)^n\) as \(n\) gets very large?

\textbf{49.} Complete the rule: When you repeatedly multiply a fraction
that is \_\_\_\_\_\_ than 1, the result approaches \_\_\_\_\_\_. When
you repeatedly multiply a fraction that is \_\_\_\_\_\_ than 1, the
result grows toward \_\_\_\_\_\_.

\textbf{50.}
\(\left(\dfrac{1}{2}\right)^1 + \left(\dfrac{1}{2}\right)^2 + \left(\dfrac{1}{2}\right)^3 + \left(\dfrac{1}{2}\right)^4\)

\textbf{51.} \(\left(\dfrac{9}{10}\right)^{10}\) --- estimate first,
then calculate. Is it closer to 0, 0.5, or 1?

\textbf{52.} \(\left(\dfrac{99}{100}\right)^{100}\) --- without
calculating exactly, estimate whether this is closer to 0, 0.3, 0.5, or
1. Explain your reasoning.

\begin{center}\rule{0.5\linewidth}{0.5pt}\end{center}

\hypertarget{part-g-basic-what-happens-when-8-problems}{%
\subsection{\texorpdfstring{Part G: Basic ``What Happens When\ldots{}''
(8
problems)}{Part G: Basic ``What Happens When\ldots'' (8 problems)}}\label{part-g-basic-what-happens-when-8-problems}}

No calculation required --- describe the behavior.

\textbf{53.} Consider \(\dfrac{n}{100}\) where \(n > 0\). - (a) As \(n\)
gets very large? - (b) As \(n\) approaches 0?

\textbf{54.} Consider \(\dfrac{100}{n}\) where \(n > 0\). - (a) As \(n\)
gets very large? - (b) As \(n\) approaches 0?

\textbf{55.} For \(\dfrac{k + 1}{k}\) where \(k > 0\): - (a) Rewrite as
a sum of two fractions. - (b) As \(k \to \infty\), what value does this
approach?

\textbf{56.} For \(\dfrac{t^2}{t^2 + 1}\) where \(t > 0\): - (a)
Calculate when \(t = 1\), \(t = 10\), \(t = 100\). - (b) What value does
this approach as \(t \to \infty\)? - (c) Will it ever actually reach
that value?

\textbf{57.} For \(\dfrac{1}{1 + n}\) where \(n > 0\): - (a) As
\(n \to 0\)? - (b) As \(n \to \infty\)?

\textbf{58.} For \(\dfrac{n}{n + 1}\) where \(n > 0\): - (a) As
\(n \to \infty\)? - (b) Is the fraction always less than 1, equal to 1,
or greater than 1?

\textbf{59.} You have fraction \(\dfrac{a}{b}\) where \(a, b > 0\). -
(a) If you double both numerator and denominator, what happens to the
value? - (b) If you add 1 to both numerator and denominator, does the
value increase, decrease, or depend on something? \emph{(Test with
\(\frac{1}{4}\) and \(\frac{3}{4}\))}

\textbf{60.} For \(\left(1 + \dfrac{1}{n}\right)\) where \(n > 0\): -
(a) Calculate when \(n = 1\), \(n = 10\), \(n = 100\). - (b) What value
does this approach as \(n \to \infty\)? - (c) Bonus:
\(\left(1 + \dfrac{1}{n}\right)^n\) approaches a famous number as
\(n \to \infty\). Any guesses?

\begin{center}\rule{0.5\linewidth}{0.5pt}\end{center}

\end{document}
